\documentclass[11pt]{article}
\usepackage[a4paper,margin=23mm]{geometry}
\usepackage{amsmath,amssymb,amsthm,mathtools,bm}
\usepackage{booktabs,longtable,array}
\usepackage{enumitem}
\usepackage{xcolor,microtype}
\usepackage[hidelinks]{hyperref}
\setlength{\emergencystretch}{3em}
\newtheorem{proposition}{Proposition}[section]
\newtheorem{lemma}[proposition]{Lemma}
\theoremstyle{definition}
\newtheorem{gate}[proposition]{Fail-closed gate}
\theoremstyle{remark}
\newtheorem{remark}[proposition]{Remark}
\newcommand{\dd}{\mathrm d}
\newcommand{\Tr}{\operatorname{Tr}}
\newcommand{\cH}{\mathcal H}
\newcommand{\RR}{\mathbb R}
\newcommand{\one}{\mathbf 1}
\newcolumntype{L}[1]{>{\raggedright\arraybackslash}p{#1}}

\title{AP-E6 Deterministic Classical Gate:\\
A Coupled Relaxed \(B=1\) Solution, Multigrid Scan,\\
and Exact Sparse \(n+s\) Projected-Curvature Audit}
\author{Route F research note}
\date{16 July 2026}

\begin{document}
\maketitle

\begin{abstract}
We replace the analytic-profile and single-field limitations of the preceding
charged two-colour proxy by a coupled relaxation.  The declared classical
sector contains a unit meson field \(n:\RR^3\to S^3\), a neutral breathing
field \(s\), a charge-two complex diquark at zero background, and free
gauge--ghost quadratic blocks.  A two-field hedgehog \((F(r),s(r))\) is solved
as a nonlinear boundary-value problem on four radial volumes.  The same
relaxed solution is sampled on three lattice spacings and three physical
volumes.  Starting from one explicitly stated finite-difference energy, we
assemble its analytic sparse embedding Hessian, apply the constrained
Riemannian correction, restrict it to \(T_nS^3\), and remove three translations
and three isorotations by an explicit orthogonal projector.  Toron and
constant-ghost projectors are audited separately.  Negative controls detect
a charged tachyon, an omitted collective direction, and loss of Derrick
stationarity without the Skyrme term.

The cubic-grid audit deliberately exposes an obstruction.  The sampled
continuum solution has tangential gradient density approximately
\(0.35\)--\(0.43\), while the displayed lattice baryon estimator only rises
from \(0.168\) to \(0.625\).  It is therefore off shell for the finite cubic
energy.  Its negative projected curvatures cannot be called physical
instability eigenvalues; instead they fail-close the discrete-stationarity
and topology-convergence gates.

Completeness below refers only to exact differentiation of the declared
deterministic classical sector.  No fermion determinant, interacting
colour-link ensemble, four-dimensional importance sampling, HMC, or quantum
continuum limit is computed.  Consequently the microscopic charged-QC2D lane
and the degree-one Route-E portal remain closed.
\end{abstract}

\tableofcontents

\section{Scope and epistemic levels}

Two-colour QCD has pseudoreal matter and genuine meson--diquark competition
\cite{KogutEtAl2000APE6}.  Recent linear and extended-linear sigma analyses
make clear that scalar and spin-one channels can reorganise in dense matter
\cite{SuenagaEtAl2023APE6,SuenagaEtAl2024APE6}.  A classical EFT soliton is
therefore neither a lattice measurement nor a proof of the microscopic phase.
Four-dimensional simulations exist on much larger ensembles
\cite{BegunEtAl2022APE6}; a recent tensor calculation controls a different
one-dimensional regulator \cite{PaiAkiyamaTodo2025APE6}.  Neither supplies the
charged scalar, soliton background, or projected Hessian studied here.

We distinguish four statements:
\begin{enumerate}[label=(L\arabic*)]
\item a stationary solution of an explicitly declared classical EFT;
\item the exact second derivative of a declared discretisation, including
  its collective-coordinate quotient, with a separate stationarity test;
\item a dynamical four-dimensional Euclidean path integral;
\item a regulator-independent quantum continuum phase.
\end{enumerate}
This note reaches L1 only in the coupled hedgehog sector; no independent
principle-of-symmetric-criticality proof is invoked to promote it to arbitrary
three-dimensional variations.  It reaches the algebraic differentiation part
of L2.  The sampled cubic configuration fails the discrete-stationarity and
topology tests, so its L2 operator is an off-shell curvature diagnostic.

\begin{gate}[Meaning of ``complete'']
Sparse differentiation is complete only for the \(n+s\) block:
every entry is a derivative of one energy, and a geodesic energy difference
checks it.  The \(\Delta\) spectrum is a separate Dirichlet block, while the
gauge/ghost algebra uses a separate periodic \(5^3\) grid.  They are not one
same-grid assembled Hessian.  Moreover, exact differentiation does not make
the sampled cubic field stationary.  Thus the aggregate
\path{sparse_projected_hessian_complete_in_declared_sector} gate is false,
and this is not yet a physical stability Hessian nor the complete
gauge--meson--ghost--fermion Hessian of charged QC2D.  In particular, there is
no Wilson/overlap fermion determinant on the same relaxed background.
\end{gate}

\section{Declared coupled classical EFT}

\subsection{Fields and energy}

Let \(n=(n_0,\bm n)\in S^3\), \(n\cdot n=1\), and let \(s\in\RR\) be a
neutral breathing field with vacuum value one.  In units \(f=e=1\), define
\begin{align}
 E[n,s]=\int_{\RR^3}\!\dd^3x\,\bigg\{&
 \frac12\partial_i n\cdot\partial_i n
 +\frac{\kappa}{2}\sum_{i<j}
 \left[|\partial_i n|^2|\partial_j n|^2
 -(\partial_i n\cdot\partial_j n)^2\right]\nonumber\\
 &+\frac12|\nabla s|^2+\frac12m_s^2(s-1)^2
 +\beta^2s^2(1-n_0)\bigg\}.
 \label{eq:energy}
\end{align}
The first line is the sigma-plus-Skyrme energy
\cite{Skyrme1961APE6,AdkinsNappiWitten1983APE6}.  The last term couples the
breathing field to explicit chiral breaking while preserving the vacuum
\((n,s)=(\one,1)\).  It is a declared diagnostic interaction, not a UV
matching result.  The benchmark is
\begin{equation}
 \beta=\frac12,\qquad m_s=\frac52,\qquad \kappa=1.
 \label{eq:parameters}
\end{equation}

A charge-two complex scalar is retained only to quadratic order about
\(\overline\Delta=0\):
\begin{equation}
 \cH_\Delta=-D_i^{(2)}D_i^{(2)}+m_\Delta^2
 +\chi(1-\overline n_0),\qquad
 m_\Delta^2=0.81,\quad\chi=-0.30.
 \label{eq:diquark}
\end{equation}
At zero charged background, gauge--diquark mixing vanishes.  This does not
cover a Higgsed/condensed branch.

\subsection{Topology and the coupled hedgehog}

Use
\begin{equation}
 n_0=\cos F(r),\qquad \bm n=\widehat{\bm x}\sin F(r),
 \qquad s=s(r),
 \label{eq:hedgehog}
\end{equation}
with \(F(0)=\pi\), \(F(\infty)=0\), \(s'(0)=0\), and
\(s(\infty)=1\).  Then
\begin{equation}
 B=-\frac{1}{2\pi^2}\int\det(n,\partial_1n,\partial_2n,
 \partial_3n)\,\dd^3x
 =-\frac2\pi\int_0^\infty\sin^2F\,F'\,\dd r=+1.
 \label{eq:baryon}
\end{equation}
Substitution into Eq.~\eqref{eq:energy} gives \(E=4\pi\int\mathcal E_r\dd r\),
where the dimensionless radial integrand is
\begin{align}
 \mathcal E_r={}&\frac12r^2F'^2+\sin^2F
 +\kappa\left(\sin^2F F'^2+\frac{\sin^4F}{2r^2}\right)
 +\frac12r^2s'^2\nonumber\\
 &+\frac12m_s^2r^2(s-1)^2+\beta^2r^2s^2(1-\cos F).
 \label{eq:radial-energy}
\end{align}

\subsection{Euler--Lagrange boundary-value problem}

Direct variation gives
\begin{align}
 (r^2+2\kappa\sin^2F)F''={}&-2rF'
 -2\sin F\cos F(\kappa F'^2-1)
 +\frac{2\kappa\sin^3F\cos F}{r^2}
 +\beta^2r^2s^2\sin F,
 \label{eq:F-bvp}\\
 s''+\frac2r s'={}&m_s^2(s-1)+2\beta^2s(1-\cos F).
 \label{eq:s-bvp}
\end{align}
Near the origin,
\begin{equation}
 F=\pi-ar+br^3+O(r^5),\qquad
 b=\frac{a(2\kappa a^4+4a^2+3\beta^2s_0^2)}
 {30(1+2\kappa a^2)},
 \label{eq:origin-F}
\end{equation}
and
\begin{equation}
 s=s_0+\frac16[m_s^2(s_0-1)+4\beta^2s_0]r^2+O(r^4).
 \label{eq:origin-s}
\end{equation}
These series, rather than a hard-coded profile, seed the nonlinear BVP.

\begin{proposition}[Generalised Derrick identity]
Every sufficiently localised stationary solution of
Eq.~\eqref{eq:radial-energy} satisfies
\begin{equation}
 E_2+E_{s,\mathrm{kin}}-E_4
 +3(E_{s,\mathrm{pot}}+E_m)=0.
 \label{eq:virial}
\end{equation}
Its scale curvature is \(2E_4+6(E_{s,\mathrm{pot}}+E_m)\).
\end{proposition}
\begin{proof}
Apply \(F(r)\mapsto F(r/\lambda)\),
\(s(r)\mapsto s(r/\lambda)\).  Two-derivative terms scale as \(\lambda\),
the Skyrme term as \(\lambda^{-1}\), and potential terms as
\(\lambda^3\).  Differentiation at \(\lambda=1\) proves
Eq.~\eqref{eq:virial}; a second derivative gives the stated curvature.
\end{proof}

\section{One lattice energy and its exact sparse Hessian}

\subsection{Declared discretisation}

On an open cubic lattice of spacing \(a\), freeze the boundary to the sampled
continuum profile rather than overwriting it by the vacuum,
let \(\delta_i n_x=(n_{x+\hat\imath}-n_x)/a\).  The single energy used by
the implementation is
\begin{align}
 E_a={}&a^3\sum_{x,i}\frac12(|\delta_i n_x|^2+|\delta_i s_x|^2)
 +\frac{\kappa a^3}{2}\sum_{x,i<j}
 \left(|\delta_i n_x|^2|\delta_jn_x|^2
 -(\delta_i n_x\cdot\delta_jn_x)^2\right)\nonumber\\
 &+a^3\sum_x\left[\frac12m_s^2(s_x-1)^2
 +\beta^2s_x^2(1-n_{0,x})\right].
 \label{eq:lattice-energy}
\end{align}
No quadratic block is inserted by hand: every meson--breathing entry below
is an analytic second derivative of Eq.~\eqref{eq:lattice-energy}.
This sampled-profile Dirichlet condition is part of the off-shell diagnostic.
For the fixed-spacing volume scan its maximum boundary \(n\)-distance from the
vacuum decreases from about \(0.350\) at \(L=4\), through \(0.120\) at \(L=6\),
to \(0.0493\) at \(L=8\); it is not silently treated as an exact vacuum wall.

For a Skyrme cell pair \(u=\delta_i n\), \(v=\delta_j n\), write
\(W=\kappa(|u|^2|v|^2-(u\cdot v)^2)/2\).  The local blocks are
\begin{align}
 W_{uu}&=\kappa(|v|^2\one-vv^T),&
 W_{vv}&=\kappa(|u|^2\one-uu^T),\nonumber\\
 W_{uv}&=\kappa[2uv^T-vu^T-(u\cdot v)\one],&
 W_{vu}&=W_{uv}^T.
 \label{eq:wedge-hessian}
\end{align}
Together with nearest-neighbour Laplacians and the local \(s\)--\(n_0\)
cross derivative, Eq.~\eqref{eq:wedge-hessian} produces a sparse symmetric
embedding Hessian \(H_{\mathrm{emb}}\).

\subsection{Constrained tangent Hessian}

At each site choose \(E_x:\RR^3\to T_{n_x}S^3\),
\begin{equation}
 E_x^TE_x=\one_3,\qquad E_x^Tn_x=0.
\end{equation}
If \(g_x=\partial E_a/\partial n_x\) and
\(\lambda_x=n_x\cdot g_x\), the constrained Riemannian Hessian is
\begin{equation}
 H_{\mathrm{tan}}
 =T^T\left[H_{\mathrm{emb}}-\bigoplus_x
 \lambda_x\one_4\right]T,\qquad
 T_x=\begin{pmatrix}E_x&0\\0&1\end{pmatrix}.
 \label{eq:constrained-H}
\end{equation}
The multiplier term is essential: simply restricting the ambient Hessian is
not the second variation on the sphere.  Equation~\eqref{eq:constrained-H}
is assembled analytically as CSR sparse data.

An independent test chooses a horizontal vector \(v\), follows the sitewise
geodesic
\[
 n_x(t)=\cos(t|v_x|)n_x+\frac{\sin(t|v_x|)}{|v_x|}v_x,
\]
and evaluates
\begin{equation}
 \frac{E_a[n(t),s+tv_s]-2E_a[n,s]+E_a[n(-t),s-tv_s]}{t^2}
 =v^TH_{\mathrm{tan}}v+O(t^2).
 \label{eq:directional-check}
\end{equation}
This checks the Weingarten/Lagrange-multiplier term as well as the ambient
second derivative.  It checks differentiation, not stationarity.

\section{Collective and gauge projectors}

\subsection{Translations and isorotations}

Let the six sampled collective vectors be
\begin{equation}
 z_i=(E^T\partial_i n,\partial_i s),\quad i=1,2,3,
 \qquad
 z_{3+a}=(E^T(0,\bm e_a\times\bm n),0),\quad a=1,2,3.
 \label{eq:collective}
\end{equation}
After a thin QR decomposition \(Z=QR\), \(Q^TQ=\one_6\), define
\begin{equation}
 P_{\mathrm{coll}}=\one-QQ^T,\qquad
 H_{\mathrm{proj}}=P_{\mathrm{coll}}H_{\mathrm{tan}}P_{\mathrm{coll}}
 \big|_{\operatorname{im}P_{\mathrm{coll}}}.
 \label{eq:collective-projector}
\end{equation}
The implementation tests \(P^2=P\), \(Q^TP=0\), eigenvector horizontality,
and an intentionally incomplete rank-five projector.  Because a finite cube
and its Dirichlet boundary lift continuum translations, raw collective
Rayleigh quotients are reported rather than declared exact zeroes.

\subsection{Gauge and ghost zero modes}

For the decoupled free Abelian block on a periodic three-dimensional audit
grid, let \(G\) be the forward gradient and \(L=G^TG\).  Covariant gauge
fixing gives
\begin{equation}
 K_A(\xi)=\one_3\otimes L+(\xi^{-1}-1)GG^T,\qquad M_{\mathrm{FP}}=L.
 \label{eq:gauge-ghost}
\end{equation}
The three constant link torons and one constant ghost are removed by
\begin{equation}
 P_A=\one-\sum_{i=1}^3t_it_i^T,\qquad
 P_c=\one-\frac{\bm1\bm1^T}{V}.
 \label{eq:gauge-projectors}
\end{equation}
Then, up to a background-independent normalization,
\begin{equation}
 \frac12\log\det{}'K_A(\xi)-\log\det{}'M_{\mathrm{FP}}
 +\frac{V-1}{2}\log\xi
\end{equation}
is \(\xi\)-independent \cite{FaddeevPopov1967APE6}.

\section{Numerical certificate}

\subsection{Coupled radial relaxation}

The final tables and exact values are machine-generated in
\path{route_f/output/ap_e6_relaxed_b1_multigrid_hessian.json}.  The BVP is
solved independently at \(R=8,10,12,16\).  The volume sequence converges to
approximately
\begin{equation}
 E/(4\pi f/e)=6.1282155,\qquad s(0)=0.9171290,\qquad B=1,
 \end{equation}
with BVP residual below \(3\times10^{-7}\).  The nonconstant \(s(r)\) and
the change between boxes demonstrate a genuine coupled relaxation rather
than evaluation of the old analytic profile.

\subsection{Multigrid, multivolume, and spectra}

The fixed-volume scan uses \(L=8\) at \(N=9,13,17\); the fixed-spacing scan
uses \((N,L)=(9,4),(13,6),(17,8)\), all with \(a=0.5\).  Every row
re-samples the same \(R=16\) relaxed continuum solution and reassembles the
exact Hessian of Eq.~\eqref{eq:lattice-energy}.  We report the lattice
topological estimator, sparse dimension and nonzero count, tangential
stationarity residual, raw low spectrum, projected spectrum, and charge-two
diquark spectrum.

\begin{table}[ht]
\centering
\small
\begin{tabular}{rrrrrr}
\toprule
\(N\)&\(a\)&\(B_a\)&
\(\|\operatorname{grad}E_a\|_{2,\mathrm{dens}}\)&
\(\lambda_{\mathrm{proj},0}\)&\(\lambda_{\Delta,0}\)\\
\midrule
9&1.000000&0.168018&0.34766&-10.79222&1.219089\\
13&0.666667&0.440012&0.43018&-23.42272&1.223574\\
17&0.500000&0.625354&0.42171&-28.44285&1.224884\\
\bottomrule
\end{tabular}
\caption{Fixed-\(L=8\) off-shell curvature scan.  The improving but still
coarse \(B_a\) and non-small gradient prohibit a stability interpretation.}
\label{tab:grid}
\end{table}

The \(N=17\) reduced block has \(13500\) variables and \(395538\) nonzero
entries.  Its symmetry residual is below \(10^{-16}\); the best relative
residual in Eq.~\eqref{eq:directional-check} is below \(2\times10^{-9}\).
Projector idempotency, horizontality, and eigenvector-overlap residuals are
respectively below \(2\times10^{-16}\), \(10^{-13}\), and
\(10^{-8}\).  Thus the negative values in Table~\ref{tab:grid} are not
caused by omitting the sphere-curvature term or by collective leakage.  Yet
they remain off-shell finite-grid curvatures, because the background has not
been re-relaxed for Eq.~\eqref{eq:lattice-energy}.

At fixed \(a=0.5\), the negative branch is volume-stable:
\(-28.41847,-28.44005,-28.44285\).  The lowest diquark box state is unbound
and moves as \(1/L^2\); a linear fit extrapolates to \(0.813709\), close to the
declared threshold \(m_\Delta^2=0.81\), with RMS residual
\(1.28\times10^{-3}\).  These are algebraic finite-grid controls.  Since
neither cubic stationarity nor \(B_a\to1\) is demonstrated, the top-level
\path{multigrid_multivolume_converged} flag is false.

For downstream same-soliton work, the public function
\path{solve_relaxed_profile(box=16.0, sample_points=4097)} returns canonical
arrays \((r,F,s)\).  The JSON records their little-endian float64 SHA-256, so a
Callias/Yukawa calculation can prove that it consumed this exact solution.

\subsection{Negative controls}

Three falsifiers are mandatory.
\begin{enumerate}[label=(N\arabic*)]
\item Replacing \(\chi=-0.30\) by \(\chi=-5\) must create a negative
  charge-two scalar eigenvalue.
\item Removing one column from \(Q\) must leave unit leakage along the omitted
  collective direction.
\item Setting \(\kappa=0\) leaves the scale derivative
  \(E_2+E_{s,\mathrm{kin}}+3(E_{s,\mathrm{pot}}+E_m)>0\), so no finite-scale
  Derrick saddle exists.
\end{enumerate}

\section{What is and is not closed}

The following finite classical statements are closed by the verifier:
\begin{enumerate}
\item a coupled, regular, finite-volume-converged \(B=+1\) BVP exists in the
declared EFT;
\item the exact analytic sparse \(n+s\) Hessian of the declared lattice energy
is symmetric and its tangent constraint term is included;
\item translation, isorotation, toron, and constant-ghost projectors satisfy
their respective algebraic tests, with the different boundary/grid choices
recorded rather than merged;
\item the geodesic directional second difference validates the analytic
Hessian, the diquark block remains positive, and all three negative controls
fire.
\end{enumerate}

The following remain false:
\begin{enumerate}
\item stationarity of the sampled field for the cubic energy,
\(B_a\to1\), and a physical projected stability spectrum; the observed
negative curvatures are consequently off-shell diagnostics;
\item one same-grid \(n+s+\Delta+\)gauge/ghost assembled sparse Hessian;
\item the full gauge--meson--ghost--fermion Hessian on one microscopic action;
\item interacting \(SU(2)_c\) links, the fermion determinant, and measure
positivity;
\item four-dimensional importance sampling or HMC;
\item a continuum quantum phase, finite-amplitude global stability, and a
microscopic derivation of \(m_s,m_\Delta,\chi\);
\item permission to construct the degree-one Route-E portal.
\end{enumerate}

\begin{gate}[Lane status]
The calculation advances the classical soliton/Hessian lane but does not
close it at the microscopic level.  The exact machine flag names are retained
in the JSON certificate.  In prose, the false gates are:
\begin{itemize}[nosep,leftmargin=2em]
\item discrete stationarity, lattice-topology convergence, and a physical
projected Hessian at a stationary discrete solution;
\item multigrid/multivolume convergence, aggregate same-grid assembly, and
the full gauge--meson--ghost--fermion Hessian;
\item dynamical four-dimensional sampling and a continuum quantum phase;
\item physics promotion and closure of this lane.
\end{itemize}
\end{gate}

\section{Next decisive computations}

\begin{enumerate}
\item Put Wilson or overlap fermions, dynamical \(SU(2)_c\) and Abelian links,
and the charged scalar on one four-dimensional Euclidean action; then sample
it rather than freezing the background.
\item Re-relax the soliton on every four-dimensional ensemble, compute the
fermion determinant curvature or stochastic estimator, and extrapolate in
\(a\), spatial volume, Euclidean time, and statistics.
\item Test finite-amplitude escape paths in the faithful target; a positive
local projected Hessian is only a necessary condition.
\end{enumerate}

\bibliographystyle{unsrt}
\bibliography{ap_e6_relaxed_b1_multigrid_hessian}

\end{document}
